\section{Análisis del ruido pasabanda}\label{sec.ruidopasa}

Consideremos el ruido obtenido a la salida de un filtro pasabanda
de recepción:
\begin{figure}[htb]
    \centering%\begin{minipage}{\textwidth}
    \bpgf{Imagenes/ruidoPasa/ruidoPasa-1}
    \begin{blockdiagram}
        \matrix[row sep=5mm, column sep=10mm]{
             \node[punto] (inicio) {}; & \node[block] (filtro) {$[f_c-B,f_c+B]$}; & \node[etiqueta] (salida) {$\quad n(t)$};\\
        };
        \draw [->] (inicio) -- (filtro);
        \draw [->] (filtro) -- (salida);
    \end{blockdiagram}

    \endpgfgraphicnamed %\end{minipage}
\end{figure}
\FloatBarrier

$n(t)$ es estacionario, con densidad espectral de potencia
$S_n(f)$ que es de tipo pasabanda, es decir $S_n(f)=0$ fuera de la
banda $[f_c-B,f_c+B]$. Dentro de la banda, consideramos cualquier
densidad espectral, en particular no tiene por qué ser constante
en la banda ni simétrica respecto de $f_c$ (el filtro puede ser no
ideal, y/o el ruido de entrada no ser blanco). Suponemos que el
ruido es de media nula.

Para estudiar el efecto de dicho ruido en la demodulaci\'on, es conveniente tener una
representaci\'on de \emph{envolvente compleja} para el ruido:
$n_{ce}(t)=n_i(t)+jn_q(t)$ tal que
$n(t)=\Re\left[n_{ce}(t)\epi\right]$.
%
%\paragraph{Preguntas:} ¿Cómo definir $n_{ce}(t)$? ¿Es
%estacionario?

%
{\bf Preguntas:} ¿Cómo definir $n_{ce}(t)$? ¿Es estacionario?

Esta pregunta es en cierto sentido la inversa de la considerada en la secci\'on anterior.
All\'i se part\'ia de un proceso estacionario de banda base, y se buscaba un proceso modulado que tambi\'en fuese
estacionario. La respuesta tiene cierto grado de dificultad matem\'atica, por lo cual la diferimos para el Ap\'endice A.
Pero el resultado es afirmativo, lo indicamos en el siguiente enunciado:


\subsection*{Teorema}
Un proceso estacionario pasabanda $n(t)$, con densidad espectral
de potencia $S_n(f)$ concentrada en la banda
$|f|\in[f_c-B,f_c+B]$, puede escribirse como
\[n(t)=n_i(t)\cos(\omega_ct)-n_q(t)\sen(\omega_ct),\]
donde $n_i$, $n_q$ son procesos estacionarios pasabajos con igual
densidad espectral (concentrada en la banda $[-B,B]$):
\begin{center}
    \col{S_{n_i}(f)=S_{n_q}(f)=\underbrace{S_n(f-f_c)u\left(-(f-f_c)\right)}_{\oldstylenums{1}}+
    \underbrace{S_n(f+f_c)u(f+f_c)}_{\oldstylenums{2}}},
\end{center}
%
%\begin{gather*}
%    S_{n_i}(f)=S_n(f-f_c)u(f_c-f) + S_n(f+f_c)u(f+f_c),
%\end{gather*}
e igual potencia media  $E[n_i^2]=E[n_q^2]=E[n^2]$.

\noindent Además, si $S_n(f)$ es
simétrica respecto de $f_c$, entonces $n_i$, $n_q$ son incoherentes
(con correlaci\'on cruzada nula).



\paragraph{Representaci\'on gr\'afica:} Para interpretar la expresi\'on para la densidad espectral de
potencia de $n_i$, $n_q$, recurrimos al siguiente bosquejo:

\begin{figure}[htb]
    \centering\begin{minipage}{\textwidth}
    \bpgf{Imagenes/ruidoPasa/ruidoPasa-4}
        \begin{tikzpicture}
        \draw [->] (0,-6.5) -- (0,1.5) node [anchor=west] {$S_n(f)$};
        \draw [->] (-5,0) -- (5,0) node [anchor=north] {$f$};
        \begin{scope}[shift={(3,0)}]
            \draw (-1,0) .. controls (-0.5,0.3) and (-0.3,0.5) .. (0,0.5);
            \draw (0,0.5) .. controls (0.3,0.5) and (0.45,0.25) .. (0.5,0);
            \draw [dotted] (0,0) node [anchor=north] {$f_c$}--(0,0.5);
        \end{scope}

        \begin{scope}[shift={(-3,0)},xscale=-1]
            \draw (-1,0) .. controls (-0.5,0.3) and (-0.3,0.5) .. (0,0.5);
            \draw (0,0.5) .. controls (0.3,0.5) and (0.45,0.25) .. (0.5,0);
            \draw [dotted] (0,0) node [anchor=north] {$-f_c$} -- (0,0.5);
        \end{scope}

        \begin{scope}[shift={(0,-2)}]
            \draw [->] (-5,0) -- (5,0) node [anchor=north] {$f$};
            \begin{scope}[xscale=-1]
            \draw (-1,0) .. controls (-0.5,0.3) and (-0.3,0.5) .. (0,0.5);
            \draw (0,0.5) .. controls (0.3,0.5) and (0.45,0.25) .. (0.5,0);
            \end{scope}
            \path (0,0.5) node [anchor=south west] {$\oldstylenums{1}: S_n(f-f_c)u(f_c-f)$};
        \end{scope}

        \begin{scope}[shift={(0,-4)}]
            \draw [->] (-5,0) -- (5,0) node [anchor=north] {$f$};
            \draw (-1,0) .. controls (-0.5,0.3) and (-0.3,0.5) .. (0,0.5);
            \draw (0,0.5) .. controls (0.3,0.5) and (0.45,0.25) .. (0.5,0);
            \path (0,0.5) node [anchor=south west] {$\oldstylenums{2}: S_n(f+f_c)u(f+f_c)$};
        \end{scope}

        \begin{scope}[shift={(0,-6)}]
            \draw [->] (-5,0) -- (5,0) node [anchor=north] {$f$};
            \draw (-1,0) .. controls (-0.5,0.3) .. (-0.5,0.3);
            \draw (-0.5,0.3) .. controls (-0.35,0.75) and (-0.15,1) .. (0,1);
            \begin{scope}[xscale=-1]
            \draw (-1,0) .. controls (-0.5,0.25) .. (-0.5,0.3);
            \draw (-0.5,0.3) .. controls (-0.35,0.75) and (-0.15,1) .. (0,1);
            \end{scope}

            \path (1,1) node {$S_{n_i}(f)$};
        \end{scope}
    \end{tikzpicture}
    \endpgfgraphicnamed \end{minipage}
\end{figure}
\FloatBarrier
 Notar que se deduce del diagrama la identidad para la potencia total:
\[\langle{n_i^2}\rangle=E[n_i^2]
=\int S_{n_i}(f)\d f= \int \text{ \oldstylenums{1} }  + \int
\text{ \oldstylenums{2} } = \int S_n(f)\d f=E[n^2].\]



\paragraph{Nota:} Dados $n_i$, $n_q$ como arriba, para cualquier $\omega'_c$ la expresión
\[n'(t)=n_i(t)\cos(\omega'_ct)-n_q(t)\sen(\omega'_ct)\]
es siempre un proceso estacionario pasabanda, en la banda
$|f|\in[f'_c-B,f'_c+B]$.

Esto se deduce de lo visto en la secci\'on anterior para el
caso de $n_i$, $n_q$ incoherentes. Es también cierto en el caso general
(remitirse al Ap\'endice A).

